% Handbuch -- Gerüst der Harwness-Vorlage (Skill latex-report).
% Für Anleitungen und Nachschlagewerke: erst das Bild im Kopf (Überblick,
% Abbildung, Glossar), dann Schritt für Schritt, dann Referenz und
% Fehlerbehebung. Alle Texte unten sind kurze Platzhalter: ersetzen, kürzen
% oder streichen. Nichts erfinden -- was der Auftrag nicht hergibt, bleibt als
% % TODO: stehen und wird in der Übergabe genannt.
\documentclass[11pt,a4paper]{scrartcl}

% ── Anpassung: von latex.template gefüllt, Vorgaben in harw-report.sty ──
\newcommand{\harwLanguage}{%%LANGUAGE%%}
\newcommand{\harwAccentHex}{%%ACCENT%%}
\newcommand{\harwWarnHex}{%%WARN%%}
\newcommand{\harwMainFont}{%%MAINFONT%%}
\newcommand{\harwSansFont}{%%SANSFONT%%}
\newcommand{\harwMonoFont}{%%MONOFONT%%}
\usepackage{harw-report}

\title{%%TITLE%%}
\subtitle{%%SUBTITLE%%}
\author{%%AUTHOR%%}
\date{%%DATE%%}

\begin{document}
\maketitle

\begin{achtung}[Wichtig vorab]
Was vor der ersten Benutzung gilt: Voraussetzungen, Grenzen oder ein
bekannter Stolperstein.
\end{achtung}

\begin{merke}[Die Idee in drei Sätzen]
Was das Beschriebene ist. Wofür man es benutzt. Wie der typische Ablauf
aussieht.
\end{merke}

\tableofcontents
\newpage

\section{Überblick}\label{sec:ueberblick}

Ein Absatz zum Zweck und zum Leserkreis. Die Abbildung zeigt die Teile und
wie sie zusammenspielen.

\begin{figure}[htbp]
\centering
\begin{tikzpicture}[node distance=1.6cm]
  \node[harwbox,fill=harwgrau] (nutzer) {Nutzerin};
  \node[harwbox,fill=harwaccent!15,right=of nutzer] (werkzeug)
    {\textbf{Werkzeug}\\Oberfläche};
  \node[harwbox,fill=harwgrau,right=of werkzeug] (daten) {Daten\\(Ablage)};
  \draw[harwpfeil] (nutzer) -- node[harwlabel,above=2pt]{bedient} (werkzeug);
  \draw[harwdoppelpfeil] (werkzeug) -- node[harwlabel,above=2pt]{liest,\\schreibt} (daten);
\end{tikzpicture}
\caption{Die Teile und ihr Zusammenspiel (Beispiel)}\label{fig:ueberblick}
\end{figure}

\subsection{Kleines Glossar}

\begin{tabularx}{\textwidth}{@{}lL@{}}
\toprule
\textbf{Begriff} & \textbf{Bedeutung} \\
\midrule
Begriff A & Eine kurze Erklärung in einem Satz. \\
Begriff B & Eine zweite Erklärung. \\
\bottomrule
\end{tabularx}

\section{Schritt für Schritt}\label{sec:schritte}

\begin{enumerate}[label=\arabic*.,itemsep=2pt]
  \item Erster Schritt: was zu tun ist und woran man erkennt, dass er
        geklappt hat.
  \item Zweiter Schritt, zum Beispiel eine Einstellung in der Datei
        \path{einstellungen/beispiel.toml}.
  \item Dritter Schritt.
\end{enumerate}

\begin{beispiel}[Beispiel: eine Einstellung]
\begin{lstlisting}
[bereich]
name = "Beispiel"
aktiv = true
\end{lstlisting}
\end{beispiel}

\section{Referenz}\label{sec:referenz}

{\small
\begin{longtable}{@{}>{\raggedright}p{3.6cm}>{\raggedright}p{2.6cm}>{\raggedright\arraybackslash}p{8.8cm}@{}}
\toprule
\textbf{Einstellung} & \textbf{Vorgabe} & \textbf{Bedeutung} \\
\midrule
\endhead
\bottomrule
\endlastfoot
\code{name} & leer & Was die Einstellung bewirkt. \\
\code{aktiv} & \code{false} & Was die Einstellung bewirkt. \\
\end{longtable}
}

\section{Fehlerbehebung}\label{sec:fehler}

\begin{description}[leftmargin=2.6cm,style=nextline,font=\sffamily\bfseries\small]
  \item[Symptom] Was die Nutzerin sieht.
  \item[Ursache] Warum es passiert.
  \item[Abhilfe] Was sie tun kann.
\end{description}

\section{Stand}\label{sec:stand}

\begin{description}[leftmargin=2.6cm,style=nextline,font=\sffamily\bfseries\small]
  \item[Gilt für] Version oder Datum, auf das sich das Handbuch bezieht.
  \item[Offen] Was noch nicht beschrieben ist.
\end{description}

\begin{thebibliography}{9}
\addcontentsline{toc}{section}{\refname}
\small
\bibitem{quelle1} Nachname, Vorname: \emph{Titel der Quelle}. Verlag, Jahr.
  \url{https://example.org/quelle}
\end{thebibliography}

\end{document}
